% !PRINTED 359
\weberpage{359}
Da $n$ teilerfremd zu $Q$ ist, so kann man diese Zahl $= nr$
setzen, worin $r$ eine zu $Q$ teilerfremde rationale Zahl ist. Demnach hat jede Zahl $\eta'$ die Eigenschaft
\weq{1}{\eta' \equiv r \pmod{Q},}
worin $r$ eine ganze rationale Zahl und die Kongruenz in dem
Sinne $n\eta' \equiv nr$ zu verstehen ist.

Die Kongruenzen (1) lassen sich multiplizieren und dividieren,
wenn man unter $r^{-1}$ die der Bedingung $r \cdot r^{-1} \equiv 1 \pmod{Q}$
genügende ganze Zahl versteht.

Umgekehrt gehört eine ganze oder gebrochene Zahl $\eta'$, die
einer Bedingung (1) genügt, der Gruppe $O'$ an. Denn ist
$n\eta' = \omega \equiv nr$ eine Zahl in $[Q]$, so ist $\eta' = \omega : n$ ein Quotient
zweier solcher Zahlen.

3. In $O'$ ist nun wieder eine Gruppe $O_0$ enthalten, die aus
allen Zahlen $\eta_0$ besteht, die der Bedingung
\weq{2}{\eta_0 \equiv 1 \pmod{Q}}
genügen, und die Zahlen von $O_0$ ergeben durch Multiplikation
und Division gleichfalls wieder Zahlen von $O_0$.

Wir haben also dreierlei Zahlgruppen, deren jede in der
vorangehenden enthalten ist:

1. $O$ \gesperrt{enthält alle gegen $Q$ fremden Zahlen in $\Omega$.}

2. $O'$ \gesperrt{enthält die gegen $Q$ fremden ganzen und gebrochenen Zahlen der Ordnung $[Q]$.}

3. $O_0$ \gesperrt{enthält die Zahlklassen in $O'$, die nach dem Modul $Q$ mit der Einheit kongruent sind.}

Die Gruppen $O$, $O'$, $O_0$ sind unendliche \gesperrt{Abelsche} Gruppen.
Nimmt man aber ein volles Restsystem rationaler, zu $Q$ teilerfremder
 Zahlen $r_1, r_2, \ldots, r_\mu$, so kann jede Zahl in $O'$ als Produkt
eines dieser $r_i$ mit einer Zahl in $O_0$ dargestellt werden, und man
kann daher nach der in Bd.\,II, § 2 gebrauchten Ausdrucksweise
setzen:
\weq{3}{O' = r_1 O_0 + r_2 O_0 + \cdots + r_\mu O_0,}
und $\mu$ ist der \gesperrt{Index} des Teilers $O_0$ von $O$: % [sic] Gedruckt ist O ohne Strich.
\weq{4}{\mu = (O', O_0),}
also eine endliche Zahl, nämlich:
\weq{5}{\mu = \varphi(Q) = Q\, \Pi\left(1 - \frac{1}{q}\right),}
worin $\varphi(Q)$ das in der Zahlentheorie gebräuchliche Zeichen für
die Anzahl der Zahlklassen nach dem Modul $Q$ mit zu $Q$ teiler-
%CONT>>
